\section{Презентационная версия прототипа}
\headingtextsep

Презентационная версия прототипа показывает общую структуру пользовательского интерфейса, состав экранов и принятый визуальный подход без детальной реализации пользовательских действий.

\textheadingsep
\subsection{Форма представления}
\headingtextsep

Прототип выполнен в виде набора отдельных HTML-страниц с минимальными стилями CSS. Такой формат позволяет представить интерфейс в наглядном виде, выдержать единый стиль экранов и получить изображения, пригодные для включения в отчет.

\textheadingsep
\subsection{Состав экранов презентационного прототипа}
\headingtextsep

Презентационная версия включает три базовых экрана, отражающих основные пользовательские сценарии:

\begin{itemize}
\item главный экран проекта;
\item экран списка задач;
\item экран базы знаний.
\end{itemize}

На главном экране показаны верхняя панель, боковая навигация, центральная рабочая область и блок краткой подсказки. Экран задач содержит область фильтрации, перечень задач и зону просмотра выбранного элемента. Экран базы знаний демонстрирует навигацию по разделам документации и область просмотра статьи. Скриншоты указанных экранов представлены в приложении А на рисунках \ref{fig:appendix-pres-overview}, \ref{fig:appendix-pres-tasks} и \ref{fig:appendix-pres-docs}.

\textheadingsep
\subsection{Требования к оформлению}
\headingtextsep

Презентационная версия выполнена в едином сдержанном стиле. Для всех экранов использована одинаковая структура шапки, повторяющиеся навигационные элементы, светлая цветовая схема и простые прямоугольные блоки, позволяющие воспринимать макет как схему интерфейса. Основная задача этого этапа заключается в демонстрации структуры экранов и состава интерфейсных зон.